\section{Introduction}

Cardiac auscultation is the oldest and cheapest cardiac screening tool in
clinical use, but interpreting heart sounds by ear is difficult and depends
heavily on the listener's training~\cite{chizner2008cardiac,leatham1975auscultation}.
The two dominant events, the first and second heart sounds (S1 and S2), are
short transients carrying most of their energy below \SI{100}{\hertz}; murmurs
are noise-like signals occupying the intervals between them and extending to
several hundred hertz~\cite{abbas2009phonocardiography}. Automatic analysis of
the phonocardiogram (PCG) is attractive wherever expert cardiology is scarce:
primary-care triage, telemedicine, and low-cost digital stethoscopes. The task
addressed here is the binary one posed by the PhysioNet/Computing in Cardiology
(CinC) 2016 Challenge --- label a recording \emph{normal} or \emph{abnormal},
where abnormal denotes a confirmed cardiac pathology rather than merely an
audible murmur~\cite{liu2016open,clifford2016classification}.

Work on this corpus is well developed. The winning entry combined a
feature-based classifier ensemble with a convolutional network on band-limited
cardiac cycles~\cite{potes2016ensemble}; a purely convolutional system on MFCC
heat maps followed~\cite{rubin2016classifying}, as did feature ensembles that
avoid segmentation altogether~\cite{zabihi2016heart}. Wavelet-packet
descriptions of the PCG form a separate line~\cite{safara2013multi,coifman1992entropy},
motivated by the fact that the informative events are transients of tens of
milliseconds, for which the local-stationarity assumption behind the short-time
Fourier transform fits poorly~\cite{mallat2008wavelet}. Reported modified
accuracies cluster around \num{0.84}--\num{0.86}.

Less developed is the question of \emph{what} such a system learns. The public
portion of CinC 2016 comprises six sub-databases collected by different groups
with different equipment, and abnormal prevalence varies across them by roughly
\SI{69}{pp}. Site identity is therefore strongly predictive of the label, and a
model that recognises the recording device or the room can score well without
learning anything about cardiac acoustics --- the failure mode documented for
radiographic COVID-19 detectors~\cite{degrave2021ai}. Here it is not
hypothetical; it is the dominant structure in the data.

This project therefore treats the aggregate metric as a claim to be tested. We
implement one recording-level pipeline and use it to compare three
representations that differ in kind --- cepstral (MFCC), spectral (log-Mel) and
perceptual wavelet packet (PWP) --- across three classifier families. Each
performance claim is paired with a diagnostic that could refute it. $k$-means
asks whether the diagnosis is geometrically present in each feature space
\emph{before any label is used}, with cluster agreement measured against both
the diagnosis and the recording site so that a null result is still informative.
Performance is broken down per sub-database, since a single bright site exposes
a shortcut. And the explainability analysis is framed as a falsifiable
prediction: because most clinically important adult murmurs are systolic, a
detector that has learned pathology should place more attribution mass in
systole than uniform temporal attention would, which can be tested against an
analytic and a permutation null.

The scope is narrow. We propose no new architecture and make no claim beyond the
split we evaluate; the official challenge test set was never released, so all
comparison with published entries is contextual. The contribution is the
controlled comparison and the diagnostics --- and two of the three expectations
we began with were contradicted by the evidence. They are reported here in the
same voice as the results that held.
